\documentclass[11pt]{article}
\usepackage[T1]{fontenc}
\usepackage{lmodern,amsmath,amssymb,amsthm}
\usepackage[margin=1in]{geometry}
\usepackage[hidelinks]{hyperref}
\newtheorem{proposition}{Proposition}
\newtheorem{corollary}{Corollary}
\newcommand{\ind}{\mathbf{1}}
\title{A two-form semiprime sieve bound and transfer tests for gap two}
\author{\small \href{mailto:admin@cosmosgeek.org}{admin@cosmosgeek.org}\qquad \href{mailto:b321014@iiit-bh.ac.in}{b321014@iiit-bh.ac.in}}
\date{Public research draft}
\begin{document}
\maketitle
\begin{abstract}
For the one-polynomial Goldston--Graham--Pintz--Y{\i}ld{\i}r{\i}m
score with two forms and level of distribution $1/2$, we give an
upper bound of $-0.01265$ times the denominator quadratic form.
The estimate is uniform in the degree and coefficients of the
polynomial and in the admissible fixed lower factor cutoff.
We also show that affine gap-two identities obtained by odd
positive rescaling form a bipartite graph. For a geometric construction
with signed circle curvatures $-p,p+2$, we classify all integral
Descartes completions of the prescribed pair and compute a gap
of quadratic size before the next possible prime tangent curvature.
\end{abstract}

\section{The two-form quadratic form}
Goldston, Graham, Pintz and Y{\i}ld{\i}r{\i}m prove results about
small gaps between products of two primes \cite{GGPY}. Their Theorem~3
gives, for an admissible pair of linear forms, infinitely many
arguments at which both values have at most two prime factors.
Their Theorems~7 and~8 evaluate the main terms for a specific
one-polynomial sieve weight. We ask whether choosing a different
polynomial in that score can force a positive leading coefficient
for two selected $E_2$ values, where an $E_2$ is a product of two
distinct primes.

Use the notation of \cite[Thms.~7--8 and \S8, (8.1)]{GGPY} in the
case $k=2$, $\nu=1$, $B=2/\vartheta=4$. Let $P$ be a real polynomial,
$F(t)=\int_0^tP(v)\,dv$, and $h=4\eta$ the fixed lower cutoff
in the integral for the smaller prime factor. The required score,
up to a positive prefactor, is
\begin{align}
J&=\tfrac12(J_1+J_2+J_3)-J_0,
&J_0&=\int_0^1(1-t)P(t)^2\,dt . \label{score}
\end{align}
Changing variables from $x$ to $t=1-x$ and combining the two
ranges of the inner integral gives
\begin{align}
J_1+J_2
 &=\int_h^1\frac{4}{y(4-y)}
       \int_0^1\bigl(F(t)-F((t-y)_+)\bigr)^2\,dt\,dy, \label{diff}\\
J_3&=\log(3)\int_0^1 F(t)^2\,dt. \label{J3}
\end{align}
Here $(t-y)_+=\max(t-y,0)$ and the coefficient in \eqref{J3} is
$\int_1^2 4/[y(4-y)]\,dy$. The following estimate also covers
the limiting integral with $h=0$.

\begin{proposition}
For $0\le h\le1$,
\[
 J\le\left(2\log(4/3)+\tfrac38\log 3-1\right)J_0.
\]
The coefficient is approximately $-0.0126562$.
\end{proposition}
\begin{proof}
For $0<y\le1$, Cauchy--Schwarz gives
\[
 \bigl(F(t)-F((t-y)_+)\bigr)^2
       \le y\int_{(t-y)_+}^t P(v)^2\,dv.
\]
For a fixed $v$, the allowable $t$ form the interval
$[v,\min(1,v+y)]$ of length at most $1-v$. Consequently
\[
 \int_0^1\bigl(F(t)-F((t-y)_+)\bigr)^2\,dt\le yJ_0 .
\]
Integrating \eqref{diff} and enlarging the outer interval to $[0,1]$
gives $J_1+J_2\le 4\log(4/3)J_0$.

For $J_3$ use $\phi(t)=t+t^2/2$. Expanding
$F'=(F/\phi)'\phi+(F/\phi)\phi'$ and integrating the cross term
by parts gives
\[
 J_0=\int_0^1(1-t)\phi(t)^2\bigl((F/\phi)'(t)\bigr)^2\,dt
      +\int_0^1 \frac{2}{1+t/2}F(t)^2\,dt .
\]
Indeed $-((1-t)\phi')'=2t$. The boundary term
$(1-t)\phi'F^2/\phi$ vanishes at $1$, and it is $O(t)$ at $0$
since $F(0)=0$. The coefficient $2/(1+t/2)$ is at least $4/3$.
It follows that $\int_0^1F^2\le 3J_0/4$. Substitution in
\eqref{score} proves the bound.
\end{proof}
Since $J_0>0$ for $P\ne0$, every such fixed polynomial gives a
strictly negative leading score. The conclusion concerns this
quadratic form at distribution level $1/2$. In particular the
prime-inclusive two-form result of \cite[Theorem~3]{GGPY} employs
a different score.

\section{An obstruction to rescaling selected forms}
Another possible transfer starts with a family of
integer affine forms $L_i(t)=a_it+b_i$, $a_i>0$. Its aim is to
attach a gap-two identity to whichever pair of indices a selection
theorem supplies.
\begin{proposition}
Join $i$ and $j$ when there exist odd positive integers $c,d$
such that $cL_i(t)-dL_j(t)=\pm2$ identically in $t$. The resulting
graph is bipartite.
\end{proposition}
\begin{proof}
Write $a_i=2^{v_i}u_i$ with $u_i$ odd. An edge requires
$ca_i=da_j$, hence $v_i=v_j$ and $M=cu_i=du_j$ is odd.
In $\mathbb Z/4\mathbb Z$ put $z_i=b_i u_i^{-1}$. The constant
coefficient identity gives
\[
 M(z_i-z_j)=cb_i-db_j=2 \pmod4.
\]
Thus $z_i-z_j=2\bmod4$. The vertices with residues $0,1$ lie
on one side and those with residues $2,3$ on the other.
\end{proof}
In particular, no triple of forms can cover all three potential
selected pairs by identities of the specified type. A theorem
giving an unspecified pair from three indices therefore requires
additional information to force such a rescaled gap.

\section{A signed Apollonian construction for the prescribed gap}
A second way to enforce a difference of two is to prescribe the
signed pair of curvatures $-p,p+2$ and apply a theorem on primes
tangent to the circle of curvature $-p$. For a growing prime anchor
$p$, this would require a tangent prime below a cutoff depending on
$p$. We compute that cutoff and the range of the fixed-form estimate.
A Descartes quadruple is a vector of oriented curvatures satisfying
\[
 2(a^2+b^2+c^2+d^2)=(a+b+c+d)^2 .
\]
Its Apollonian packing is obtained by the integral Descartes
reflections. A circle's tangent curvatures admit a positive-definite
shifted binary quadratic form parametrization due to Sarnak; we use
Proposition~1.2 of \cite{FHRS}. In a quadruple $(a,b,c,d)$, primitive
pairs $(x,y)$ give the tangent curvatures at $a$ as the values of
\[
 f_a(x,y)-a=(b+a)x^2+(a+b+d-c)xy+(d+a)y^2-a .
\]

\begin{proposition}
Fix a prime $p\ge5$. Write $T=(p+1)^2/2$.
All ordered integral completions $(-p,p+2,c,d)$ have the form
\[
             c=T+2h(h-1),\qquad d=T+2h(h+1),\qquad h\in\mathbb Z .
\]
The resulting unshifted tangent forms at $-p$ belong to one proper
equivalence class of discriminant $-4p^2$. In each case $p+2$ occurs
as a tangent curvature, and every prime tangent curvature other than
$p+2$ is at least $2p^2+p+8$.
If $p\equiv2\bmod3$, every such prime is at least $2p^2+p+72$.
\end{proposition}
\begin{proof}
The Descartes equation rearranges to
$(a-b)^2+(c-d)^2=2(a+b)(c+d)$.
For $a=-p,b=p+2$ this becomes
\[
              4(p+1)^2+(d-c)^2=4(c+d).
\]
Thus $d-c$ is even. Since $c+d$ is then even, reduction modulo eight
gives $4\mid d-c$. Set $d-c=4h$; substitution and taking half the sum
and difference gives the stated formulas. Conversely these formulas
satisfy the equation. The quadruple is primitive since $(p,p+2)=1$,
and $h=0$ gives $(-p,p+2,T,T)$.

Substitution in the tangent form gives
\[
 f_h(x,y)=2x^2+(2+4h)xy+
 \left((p^2+1)/2+2h(h+1)\right)y^2=f_0(x+hy,y).
\]
The change of variables has determinant one and preserves primitivity.
The form $f_0$ is primitive because $(p^2+1)/2$ is odd, and its
discriminant is $-4p^2$. It suffices to examine
\[
 q_p(x,y)=p+f_0(x,y)=p+\tfrac12\bigl((2x+y)^2+p^2y^2\bigr).
\]
If $y=0$ then $x=\pm1$, yielding $p+2$. If $y$ is odd then the
two squares in parentheses are $1\bmod8$, so $q_p$ is even and
greater than $2$. Otherwise write $y=2z$, where $z\ne0$ and $x$ is
odd. Now
\[
 q_p(x,2z)=p+2(x+z)^2+2p^2z^2.
\]
For $|z|=1$ the integer $x+z$ is even. When $x+z=0$ the value is
$p(1+2p)$; otherwise $|x+z|\ge2$. For $|z|\ge2$ the value is at
least $p+8p^2$. This proves the first bound. If $p\equiv2\bmod3$
and $|z|=1$, the value is $1+2(x+z)^2\bmod3$. A prime in this
range is greater than three, forcing $3\mid x+z$ and hence
$|x+z|\ge6$ unless $x+z=0$. For $|z|\ge2$ the preceding bound
already exceeds $2p^2+p+72$.
\end{proof}

For each fixed $p$, Theorem~1.3 of \cite{FHRS} gives prime
curvature counts of order $B/(\log B)^{3/2}$ as $B$ grows.
In contrast, for $B=Cp^2$ the display for $q_p(x,2z)$ shows
$|z|\le\sqrt{C/2}$ and $|x+z|\le p\sqrt{C/2}$. There are at
most $O_C(p)$ pairs and hence that many distinct prime curvatures,
including the possible value $p+2$. A positive lower constant
independent of $p$ at this scale would contradict this count.
Changing the completion $(c,d)$ leaves all represented values unchanged.

\subsection*{Direct congruence restrictions}
\begin{proposition}
All odd curvatures in a primitive integral Apollonian packing are
congruent modulo four. Moreover, if $p,p+2$ are primes with $p>3$,
no integral Descartes quadruple contains $2p$ and $2(p+2)$.
\end{proposition}
\begin{proof}
The Descartes equation forces an even number of odd coordinates.
Primitivity excludes zero such coordinates and four odd coordinates
would require a square equal to $8\bmod16$. Write the two odd
coordinates as $a,b$. If the even coordinates $c,d$ have different
residues modulo four, the same equation would require a square equal
to $12\bmod16$. Thus $c\equiv d\bmod4$ and $a+b\equiv2\bmod4$,
giving $a\equiv b\bmod4$. Replacing an odd coordinate by twice the
sum of the others minus that coordinate preserves its residue
modulo four; the same holds for either even coordinate. All
reflections therefore preserve the odd residue class.

For $c=2p,d=2(p+2)$ the rearranged equation reads
$(a-b)^2+16=8(p+1)(a+b)$. The twin-prime condition gives
$p\equiv2\bmod3$, so the equation requires a square equal to $2$
modulo three. This is impossible.
\end{proof}


\paragraph{Provenance.}
The two-form estimates and the signed Apollonian calculations were developed during the Twin Prime research project. The latter occur in replies 4 and 5 of the archived conversation \emph{Twin prime proof planning}. The derivation and drafting
involved ChatGPT and Codex. The original account calculation is
linked in the \href{https://github.com/cosmosgeek0/Twin-Prime/blob/main/history/hosted-research-responses-2026-11-25.md}{public research record}.
Specialist review and priority assessment remain open.
\begin{thebibliography}{9}
\bibitem{GGPY} D.~A.~Goldston, S.~W.~Graham, J.~Pintz and C.~Y.~Y{\i}ld{\i}r{\i}m,
\emph{Small gaps between products of two primes}, Proc. Lond. Math.
Soc. (3) \textbf{98} (2009), 741--774.
\href{https://doi.org/10.1112/plms/pdn046}{Published paper};
\href{https://arxiv.org/abs/math/0609615}{authors' preprint}.
See also the \href{https://doi.org/10.1112/plms/pds053}{2013 corrigendum}.
\bibitem{FHRS} E.~Fuchs, C.~Hsu, J.~Rickards, D.~Schindler and
K.~E.~Stange, \emph{Primes represented by shifted quadratic forms:
on primitivity and congruence classes}, Proposition~1.2 and Theorem~1.3.
\href{https://arxiv.org/abs/2504.20289}{Authors' paper}.

\end{thebibliography}
\end{document}
